\chapter[Correspondència amb les categories]{Correspondència amb el volum de categories}\label{cap:correspondencia}

Aquest llibre no fa servir el vocabulari de la teoria de categories, però
cada construcció seva és el cas particular, en els conjunts, d'una noció
categòrica general. La taula següent indica, per a cada construcció, el nom
categòric que rep i el capítol del volum \emph{Una introducció a la teoria de
categories}~\cite{perello-categories} on es tracta en general. Els capítols
d'aquell volum són: 1~Categories, 2~Functors, 3~Transformacions naturals,
4~Yoneda i propietats universals, 5~Límits i colímits, 6~Adjuncions,
7~Categories cartesianes tancades, 8~Subobjectes i categories regulars,
9~Topos elementals, i 10~el pont cap a la lògica categòrica.

{\small
\begin{longtable}{@{}>{\raggedright\arraybackslash\hspace{0pt}}p{0.34\linewidth}>{\raggedright\arraybackslash\hspace{0pt}}p{0.36\linewidth}>{\raggedright\arraybackslash\hspace{0pt}}p{0.20\linewidth}@{}}
\toprule
Aquest llibre & Noció categòrica & Categories\\
\midrule
\endfirsthead
\toprule
Aquest llibre & Noció categòrica & Categories\\
\midrule
\endhead
\bottomrule
\endfoot
Aplicacions, composició, identitats (cap.~\ref{cap:aplicacions}) & Categoria; la dels conjunts, \(\mathbf{Set}\) & cap.~1\\
Isomorfisme; dualitat (cap.~\ref{cap:aplicacions}) & Isomorfisme; categoria oposada & cap.~1\\
Element generalitzat \(X\to A\) & Morfisme amb codomini \(A\); punt de vista de Yoneda & cap.~4\\
Terminal i inicial (cap.~\ref{cap:aplicacions}) & Objectes terminal i inicial & cap.~4\\
Separació global & Categoria ben apuntada; \(\mathbf1\) separador & cap.~9\\
Monomorfisme, epimorfisme, secció, retracció (cap.~\ref{cap:parts}) & Les mateixes nocions en una categoria qualsevol & cap.~1\\
Part d'un conjunt; inclusió (cap.~\ref{cap:parts}) & Subobjecte; ordre de subobjectes & cap.~8\\
Suma i producte (cap.~\ref{cap:parts}) & Coproducte i producte & cap.~5\\
Equalitzador, pullback, imatge inversa (cap.~\ref{cap:parts}) & Límits finits; reindexació de subobjectes & caps.~5 i~8\\
Representació de les parts per \(\mathbf2\) (cap.~\ref{cap:parts}) & Classificador de subobjectes \(\Omega\) & cap.~9\\
Sumes disjuntes i estables (cap.~\ref{cap:valors}) & Coproductes disjunts i universals & caps.~5 i~9\\
Coequalitzador, pushout, límits i colímits finits (cap.~\ref{cap:valors}) & Colímits i límits finits & cap.~5\\
Relació d'equivalència; parell nucli; quocient (cap.~\ref{cap:relacions}) & Relació d'equivalència interna; \emph{kernel pair}; coequalitzador & caps.~5 i~8\\
Efectivitat de les relacions d'equivalència & Categoria exacta (de Barr) & cap.~8\\
Elecció: tot epi té secció & Axioma d'elecció interna; epis escindits & cap.~9\\
Factorització imatge (cap.~\ref{cap:imatges}) & Factorització epi regular\,--\,mono & cap.~8\\
Estabilitat dels epis; teoria regular & Categoria regular & cap.~8\\
Imatge directa \(\dashv\) imatge inversa (cap.~\ref{cap:imatges}) & Existencial com a adjunt de la reindexació & caps.~6 i~8\\
Exponencial \(B^A\); transposició (cap.~\ref{cap:exponencials}) & Objecte exponencial; categoria cartesiana tancada; currificació & caps.~6 i~7\\
Conjunt de parts \(P(A)=\mathbf2^A\) (cap.~\ref{cap:parts-interna}) & Objecte de parts \(\Omega^A\) & cap.~9\\
Quantificadors \(\exists_f\dashv P(f)\dashv\forall_f\); Beck--Chevalley & Hiperdoctrina de subobjectes; condició de Beck--Chevalley & caps.~8 i~9\\
Objecte dels nombres naturals (cap.~\ref{cap:naturals}) & Objecte de nombres naturals (NNO) & cap.~9\\
Conjunts finits; \(\operatorname{Card}(A)\) (cap.~\ref{cap:finits}) & Objectes finits (Kuratowski) en un topos & ---\\
Famílies \(E\to I\); fibres (cap.~\ref{cap:families}) & Categoria de fletxes \(\mathbf{Set}/I\) & caps.~4 i~9\\
Bons ordres, ordinals, cardinals (caps.~\ref{cap:bons-ordres}--\ref{cap:cardinals}) & Sense contrapartida directa al volum de categories & ---\\
Esquema SR (cap.~\ref{cap:alephs}) & Reemplaçament estructural & ---\\
Tot el sistema d'axiomes (cap.~\ref{cap:models}) & Topos ben apuntat amb NNO i elecció: ETCS & caps.~9 i~10\\
\end{longtable}
}

Dues lectures de la taula. La primera: tot el que el llibre fa fins al
capítol~\ref{cap:parts-interna} és teoria de topos aplicada a un topos
concret; el volum de categories ho fa per a un topos qualsevol, on ja no hi
ha prou punts ni elecció, i on el classificador pot tenir més de dos valors.
La segona: les files sense contrapartida ---la teoria transfinita i SR--- són
el que és propi de la teoria de \emph{conjunts}, i no de les categories: la
mida i la longitud, que només tenen sentit quan es pot comparar i ordenar
sense límit.
